% !TEX root = ../algorithms.tex

\section{Хэш-таблица с открытой адресацией}

\subsection{Основная идея и структура данных}

\begin{Definition}{}{}
	Хранится таблица $T[0..m-1]$ такая, что $n < m$, где $n$ — количество ключей, хранящихся в таблице в данный момент. Зафиксируем функцию $h\colon U \to \{0, \ldots, m-1\}$, которая задаёт ячейки, в которые помещаются ключи: $x \in U$ помещается в $T[h(x)]$. Функция $h$ называется \emph{хэш-функцией}. Универсум ключей $U = \{0, 1, \ldots, m-1\}$.
\end{Definition}

\begin{Definition}{}{}
	\emph{Коллизия} — это ситуация, когда $h(x) = h(y)$ при $x \neq y$.
\end{Definition}

\begin{Statement}{}{}
	Если $|U| > m$, то коллизии неизбежны при любой хэш-функции $h$.
\end{Statement}

\subsection{Разрешение коллизий с помощью линейного поиска}

При линейном зондировании ячейки таблицы просматриваются последовательно, начиная с позиции $h(x)$, по кругу. Ниже приведена иллюстрация: стрелка указывает на позицию $h(x)$, серые ячейки образуют \emph{серию} (цепочку подряд идущих занятых ячеек).

\begin{center}
	\begin{tikzpicture}[scale=0.85]
		% Draw circular hash table
		\def\R{2.2}
		\def\N{18}
		% Draw empty cells
		\foreach \i in {0,...,17} {
			\pgfmathsetmacro{\angle}{90 - \i * 360/\N}
			\pgfmathsetmacro{\ax}{\R * cos(\angle)}
			\pgfmathsetmacro{\ay}{\R * sin(\angle)}
			\node[draw, circle, inner sep=2pt, minimum size=10pt,
			fill=white] at (\ax, \ay) {};
		}
		% Mark occupied cells (a run/series), e.g. positions 2..6
		\foreach \i in {2,3,4,5,6} {
			\pgfmathsetmacro{\angle}{90 - \i * 360/\N}
			\pgfmathsetmacro{\ax}{\R * cos(\angle)}
			\pgfmathsetmacro{\ay}{\R * sin(\angle)}
			\node[draw, circle, inner sep=2pt, minimum size=10pt,
			fill=gray!60] at (\ax, \ay) {};
		}
		% Arrow pointing to h(x) = position 4
		\pgfmathsetmacro{\hangle}{90 - 4 * 360/\N}
		\pgfmathsetmacro{\hx}{\R * cos(\hangle)}
		\pgfmathsetmacro{\hy}{\R * sin(\hangle)}
		\draw[->, thick] (\hx + 0.9, \hy + 0.9) -- (\hx + 0.18, \hy + 0.18);
		\node at (\hx + 1.2, \hy + 1.05) {$h(x)$};
		% Label "серия"
		\node at (4.2, -0.3) {серия};
		\draw[->, black!70] (3.5, -0.2) -- (2.6, -0.35);
	\end{tikzpicture}
\end{center}

\subsection{Алгоритмы операций}

В таблице используется специальный маркер $\dagger$ (tombstone) для обозначения удалённых ячеек, а символ $\varnothing$ — для обозначения пустой ячейки.

\subsubsection*{Вставка}

\begin{algorithmic}[1]
	\Function{insert}{$x, v$}
	\State $t = h(x)$
	\While{$T[t] \neq \varnothing$ \textbf{and} $T[t] \neq \dagger$ \textbf{and} $T[t].\mathit{key} \neq x$}
	\State $t = (t + 1) \bmod m$
	\EndWhile
	\If{$T[t].\mathit{key} \neq x$}
	\State $T[t].\mathit{key} = x$
	\State $n = n + 1$
	\EndIf
	\State $T[t].\mathit{value} = v$
	\If{$2n > m$}
	\State Создаём новую пустую таблицу $T'$ размера $2m$
	\State Перехэшируем все элементы из $T$ в $T'$ с учётом нового размера таблицы
	\State Удаляем $T$
	\State $T = T'$
	\State $m = 2m$
	\EndIf
\end{algorithmic}

\subsubsection*{Удаление}

\begin{algorithmic}[1]
	\Function{delete}{$x$}
	\State $t = h(x)$
	\While{$T[t] \neq \varnothing$ \textbf{and} $T[t].\mathit{key} \neq x$}
	\State $t = (t + 1) \bmod m$
	\EndWhile
	\If{$T[t] \neq \varnothing$}
	\State $T[t] = \dagger$
	\EndIf
\end{algorithmic}

\subsection{Идеализированное предположение и анализ}

\begin{Definition}{Идеализированное предположение}{}
	Хэш-функция $h$ генерируется случайно и равновероятно из множества всех функций $U \to \{0, \ldots, m-1\}$, то есть для каждого $x \in U$ значение $h(x)$ выбирается независимо и равновероятно из $\{0, \ldots, m-1\}$.
\end{Definition}

\begin{Theorem}{}{}
	При идеализированном предположении ожидаемое время работы каждой операции равно $O(1)$ (если не происходит перестройка таблицы).
\end{Theorem}

\begin{proof}
	Будем считать, что все ключи различны и $m = 2n$.
	
	Время работы операции пропорционально длине «хвоста» серии, начиная с позиции $h(x)$. Обозначим:
	\begin{itemize}
		\item \emph{серия} — максимальная цепочка подряд занятых ячеек $t', t'+1, \ldots, t'+k-1$, где $t'-1$ и $t'+k$ пусты (нумерация по модулю $m$);
		\item $q_k$ — вероятность того, что в таблице существует серия длины $k$, начинающаяся в фиксированной позиции.
	\end{itemize}
	
	\noindent\textbf{Ожидаемое время.}
	Пусть $p_\ell$ — вероятность встретить в позиции $t = h(x)$ «хвост» длины $\ell$ некоторой серии длины $k \geq \ell$.
	Если серия имеет длину $k \geq \ell$, то она должна начинаться в позиции $t - (k - \ell)$, поэтому
	\[
	p_\ell = \sum_{k=\ell}^{n} q_k.
	\]
	Тогда ожидаемое время равно
	\[
	O\!\left(\sum_{\ell=0}^{n} p_\ell \cdot \ell\right).
	\]
	Преобразуем:
	\begin{align*}
		\sum_{\ell=0}^{n} p_\ell \cdot \ell
		&= \sum_{\ell=0}^{n} \sum_{k=\ell}^{n} q_k \cdot \ell
		= \sum_{k=0}^{n} q_k \cdot (0 + 1 + 2 + \cdots + k)
		= \sum_{k=0}^{n} q_k \cdot \Theta(k^2).
	\end{align*}
	
	\noindent\textbf{Оценка $q_k$.}
	Докажем, что $q_k \leq e^{-\Omega(k)}$. Из этого будет следовать:
	\begin{align*}
		\sum_{k=0}^{n} q_k \cdot \Theta(k^2)
		\leq \Theta\!\left(\sum_{k=0}^{n} \frac{k^2}{e^{\Omega(k)}}\right) = O(1).
	\end{align*}
	
	Если в таблице есть серия в позициях $t', t'+1, \ldots, t'+k-1$, значит, $k$ ключей отображены в ячейки
	$\{t', t'+1, \ldots, t'+k-1\}$, а остальные $n - k$ ключей — вне этих ячеек.
	
	Следовательно, так как $m = 2n$:
	\begin{align*}
		q_k
		&\leq \binom{n}{k}
		\left(\frac{k}{2n}\right)^k
		\left(1 - \frac{k}{2n}\right)^{n-k}.
	\end{align*}
	
	Оценим сверху каждый множитель. Используем неравенства
	\[
	\binom{n}{k} \leq \frac{n^n}{k^k (n-k)^{n-k}},
	\qquad
	1+x \leq e^x.
	\]
	
	Первое неравенство следует из бинома Ньютона: $n^n = (k + (n-k))^n
	= \sum_{i=0}^{n} \binom{n}{i} k^i (n-k)^{n-i}.$
	
	В этой сумме все слагаемые неотрицательны, поэтому $n^n \geq \binom{n}{k} k^k (n-k)^{n-k}.$
	
	
	Тогда
	\begin{align*}
		q_k
		&\leq
		\frac{n^n}{k^k (n-k)^{n-k}}
		\cdot
		\frac{k^k}{2^k n^k}
		\cdot
		\left(\frac{2n-k}{2n}\right)^{n-k}
		\\[4pt]
		&=
		\frac{1}{2^k}
		\cdot
		\left(\frac{2n-k}{2n-2k}\right)^{n-k}
		\\[4pt]
		&=
		\frac{1}{2^k}
		\cdot
		\left(1 + \frac{k}{2n-2k}\right)^{n-k}
		\\[4pt]
		&\leq
		\frac{1}{2^k}
		\cdot
		e^{\frac{k(n-k)}{2n-2k}}
		\\[4pt]
		&=
		\frac{1}{2^k}
		\cdot
		e^{k/2}
		=
		\left(\frac{\sqrt{e}}{2}\right)^k
		=
		e^{-\Omega(k)}.
	\end{align*}
	
	Таким образом, ожидаемое время работы равно $O(1)$.
\end{proof}
